\section{Droplet velocities and conditional reference comparisons}
The PZT atomiser operated at 190 kHz and 9 Vpp. Final calibrated correlation-derived
droplet velocities were used unchanged: $5.35\,\mu\mathrm{m/pixel}$,
$\Delta t=25\,\mu\mathrm{s}$, and $0.214\,\mathrm{m\,s^{-1}}$ per pixel
displacement, applied once in the experimental analysis. Coordinates $x,u$ are
horizontal/rightward and $y,v$ upward. The measured profile is the ensemble mean
$\overline{v}(x_*,y)$ along the existing fixed line $x_*=14.54665\,\mathrm{mm}$;
it is neither an individual droplet trajectory nor a tracked plume centreline.
We use $h=y-y_0$, with image-reference $y_0=0.0856\,\mathrm{mm}$
at the first finite paired measurement used for initialisation. It is not an
established atomiser surface. The initial measured axial velocity is
$v_0=1.65501079\,\mathrm{m\,s^{-1}}$; the measured horizontal
mean is also retained in the two-component reference calculation.

Individual spherical water-droplet references are calculated in assumed stationary
air, with constant assumed diameters 5, 10, 20, 50 and $100\,\mu\mathrm{m}$,
neglecting evaporation. The air temperature is $20\,^{\circ}\mathrm{C}$;
pressure 101.325 kPa and equal water/air temperature are assumptions. Humidity is
unavailable; dry-air properties are an approximation. Adopted values are
$\rho_a=1.20479\,\mathrm{kg\,m^{-3}}$,
$\rho_w=998.20675\,\mathrm{kg\,m^{-3}}$ and
$\mu_a=1.813322\times10^{-5}\,\mathrm{Pa\,s}$.
Water density follows NIST GLP10, air density uses a NIST dry-air reference,
and viscosity uses the NASA-reported Sutherland constants
\cite{pzt-water,pzt-air,pzt-viscosity}.
\begin{align}
\dot{\boldsymbol q}&=\boldsymbol w_d,\\
\dot{\boldsymbol w}_d&=-g(1-\rho_a/\rho_w)\boldsymbol e_y
+\frac{f(Re_d)}{\tau_{p,\mathrm{Stokes}}}
(\boldsymbol w_a-\boldsymbol w_d),\qquad \boldsymbol w_a=\boldsymbol 0,\\
Re_d&=\frac{\rho_a d\lVert\boldsymbol w_a-\boldsymbol w_d\rVert}{\mu_a},
\qquad\tau_{p,\mathrm{Stokes}}=\frac{\rho_w d^2}{18\mu_a}.
\end{align}
Here $g=9.80665\,\mathrm{m\,s^{-2}}$ and the spherical drag multiplier is
$f=1$ for $Re_d\leq0.01$ and $f=1+0.15Re_d^{0.687}$ otherwise, with
$Re_d<800$ enforced \cite{pzt-drag}. Only in-plane slip is available. Rigid-interface spherical
drag, fixed diameter and one-way coupling are model assumptions; deformation,
internal circulation, evaporation and particle interactions are omitted.
The gravity-only reference omits drag and buoyancy. Each displayed upward
trajectory terminates at its turning point or the measurement boundary.

For assumed diameters 5, 10, 20, 50 and $100\,\mu\mathrm{m}$, the stationary-air
turning heights are 0.11913, 0.45743, 1.6978, 8.5551, 24.344$\,\mathrm{mm}$, respectively.
The gravity-only turning height is $139.653\,\mathrm{mm}$,
whereas the plotted measurement extent is only $27.2208\,\mathrm{mm}$.
These are different quantities. The full and $0$--$2\,\mathrm{mm}$ views
show the same unsmoothed measurements and unchanged reference curves.
Disagreement with stationary-air trajectories does not itself identify a coding
error or determine droplet size. Possible contributors include surrounding gas
motion, size and initial-velocity distributions, and lateral plume displacement;
the available evidence does not establish which dominates. Neither diameter nor
an unknown gas field is fitted, and the reference model is not experimentally validated.

Two Reynolds-number quantities are reported separately. The conditional estimate
$Re_d^{\mathrm{mean}}=\rho_a d\lVert\overline{\boldsymbol w}_d\rVert/\mu_a$
uses the measured mean vector and assumed stationary air; it is not
$\overline{Re_d}$ from instantaneous samples and spans the measurement domain.
The trajectory reference $Re_d^{\mathrm{ref}}$ uses the model's instantaneous
velocity and ends at that trajectory's endpoint. They are not equivalent results.
The response-time plot reports the Stokes-limit $\tau_{p,\mathrm{Stokes}}$.
At finite Reynolds number the exported $\tau_{\mathrm{drag,sec}}=
\tau_{p,\mathrm{Stokes}}/f(Re_d)$ is an instantaneous drag/slip coefficient
timescale, not an exact or linearised relaxation time.

The separate hypothetical jet-interaction comparison uses
$St_D(x)=\tau_{p,\mathrm{Stokes}}U_c(x)/D_0$, $D_0=2\,\mathrm{mm}$,
where $U_c(x)$ is from the separately measured ZNMF jet and its own downstream
coordinate. The assumed diameters and the local timescale $D_0/U_c(x)$ define a
local nozzle-length-based response comparison; the downstream trend reflects
that chosen timescale. This is neither a measured PZT-plume Stokes number nor
a nozzle-exit Stokes number. No spatial registration to the PZT field is assumed,
the separate jet velocity is not inserted into the PZT droplet slip, and the
calculation does not demonstrate entrainment. The 190 kHz atomiser frequency is
not a gas-flow frequency in this analysis.
